\documentclass[11pt]{article}
\usepackage[margin=1in]{geometry}
\usepackage{amsmath,amssymb,amsthm}
\usepackage[hidelinks]{hyperref}
\newtheorem{theorem}{Theorem}
\newtheorem{lemma}[theorem]{Lemma}
\newtheorem{proposition}[theorem]{Proposition}
\newtheorem{corollary}[theorem]{Corollary}
\theoremstyle{remark}
\newtheorem{remark}[theorem]{Remark}
\newcommand{\Fp}{\mathbb F_p}
\newcommand{\tr}{\operatorname{Tr}}
\newcommand{\Span}{\operatorname{span}}
\newcommand{\BCH}{\operatorname{BCH}}
\newcommand{\eps}{\varepsilon}
\title{Positive Rees Covers and Hidden Subgroup Recovery\\
\large A Self-Contained Proof Draft for Adversarial Review}
\author{Quantum Algorithm Search project\\Automated research draft; not independently reviewed}
\date{22 October 2026}
\begin{document}
\maketitle
\begin{abstract}
We describe a proposed composition of standard Lie-theoretic and quantum
algorithmic ingredients. A positive lower-central Rees cover makes an
arbitrary hidden subgroup of a prime-field Lazard group invariant under
public dilations. Central character measurements and diagonal-equation
witnesses then remove successive grades without removing the original
hidden Lie subspace from grade one. The written argument gives a
bounded-error polynomial-time procedure at every fixed nilpotency class
below the characteristic, with growing prime and dimension. A separate
corollary uses published reductions for general fixed-class nilpotent
black-box groups. For explicit Lie input, a further proposed corollary gives
uniform fixed-parameter time in the derived-algebra dimension, allowing
restricted growing class. This manuscript is a proposed proof, not independent
verification, a novelty certificate, or an implemented quantum algorithm.
It has no established lattice, graph-isomorphism, or cryptanalytic application.
\end{abstract}

\section{Claim, status, and input model}
Fix a constant $c$. Let $p>c$ be an odd prime, and let $L$ be an
$n$-dimensional nilpotent Lie algebra over $\Fp$ of class at most $c$,
specified by genuine structure constants in a public basis. Its Lazard
group $G=\exp(L)$ has multiplication given by the finite BCH formula.
The correspondence preserves subgroups and Lie subalgebras
\cite[Section 2.3]{lazard}; in particular $\mathfrak h=\log H$ is an
$\Fp$-linear subspace for every $H\le G$.

The ordinary input is a coherent fixed-width function oracle
$U_f|g,z\rangle=|g,z\mathbin{\mathrm{XOR}} f(g)\rangle$ satisfying
$f(g)=f(g')$ exactly when $g^{-1}g'\in H$. All oracle calls, label bits,
public arithmetic, reversible cleanup and approximation costs are charged.
Alternatively, the input can be fresh IID copies of
\[
 \rho_H=\frac1{|G|}\sum_{h\in H}R_h,
 \qquad R_g|x\rangle=|xg^{-1}\rangle,
\]
together with the same public group operations. No inverse of the unknown
source, hidden coefficient, membership oracle for $H$, or conjugate copy
is assumed. A correlated bank with these one-copy marginals is insufficient.
There is a sufficient weaker raw-input promise: an arbitrary mixture of
products of ideal coset states can be refreshed into independent samples,
as shown below. Separability alone is not this promise.

\begin{theorem}[Proposed explicit-Lie theorem]\label{thm:main}
Under the preceding input contract, the algorithm below returns a basis
of $\mathfrak h$ with probability at least $1-\eps$, in time polynomial
in $n$, $\log p$, the supplied oracle-label width, and $\log(1/\eps)$,
for fixed $c$. The only classical non-elementary arithmetic dependency
is the published fixed-degree diagonal-system algorithm of
Ivanyos--Santha \cite{is}. The statement is bounded-error, not exact
over a fixed finite quantum gate set.
\end{theorem}

The proofs below are the project's proposed justification of this theorem.
Finite tests are not premises of the theorem and do not certify it. The
main new item to scrutinize is the composition of homogenization and
all-record central descent, not the standard ingredients separately.

\section{A public homogeneous cover}
Write $\gamma_1(L)=L$ and $\gamma_{r+1}(L)=[L,\gamma_r(L)]$. Set
\[
 A=\bigoplus_{r=1}^{c} A_r,\qquad
 A_r=t^r\gamma_r(L)\subset L[t],\qquad
 D=\dim A\le cn.
\]
The bracket is the actual polynomial bracket. Since
$[\gamma_r(L),\gamma_s(L)]\subseteq\gamma_{r+s}(L)$ and
$\gamma_{c+1}(L)=0$, this is a Lie subalgebra; no nonzero bracket is
silently truncated. It has class at most $c$. Evaluation
$E_s(a(t))=a(s)$ is a Lie homomorphism for every $s\in\Fp$.
For $s\ne0$, $\Delta_s(a(t))=a(st)$ is an automorphism scaling
grade $r$ by $s^r$. These maps are also BCH group homomorphisms.

Choose distinct nonzero $s_1,\ldots,s_c$, for instance $1,\ldots,c$.
Define the tuple oracle on $\exp(A)$ by
\[
 \widehat f(a)=\bigl(f(\exp E_{s_1}(\log a)),\ldots,
                         f(\exp E_{s_c}(\log a))\bigr).
\]
It uses $c$ ordinary oracle calls. Each evaluation workspace is uncomputed
while retaining the address and the tuple label. Discarding uncleaned
evaluation registers would generally destroy the required coherence.

\begin{lemma}[Exact homogenization]\label{lem:homogeneous}
The tuple oracle hides $\widetilde H=\exp(\widetilde{\mathfrak h})$, where
\[
 \widetilde{\mathfrak h}
 =\bigoplus_{r=1}^{c}t^r(\mathfrak h\cap\gamma_r(L)).
\]
This subgroup is preserved by every nonzero $\Delta_s$.
\end{lemma}
\begin{proof}
For addresses $a,b$, write $\log(a^{-1}b)=\sum_r C_rt^r$.
Functoriality of BCH gives $\widehat f(a)=\widehat f(b)$ precisely when
$\sum_r s_i^r C_r\in\mathfrak h$ for each $i$.
The matrix $(s_i^r)_{i,r=1}^c$ has determinant
$\prod_i s_i\prod_{i<j}(s_j-s_i)\ne0$.
Linear combinations using its inverse show this is equivalent to
$C_r\in\mathfrak h$ for every $r$. The displayed subspace is a Lie
subalgebra and is homogeneous, giving the final assertion.
\end{proof}

In particular $A_1=tL$, \emph{not} $L/[L,L]$. Duplicating derived directions
in grade one is essential. In the associated graded algebra they would
already have been lost before the quantum procedure began.

\begin{remark}[No original positive grading is assumed]
The known eight-dimensional Dixmier--Lister algebra
\cite[Section 4.2, Example 2]{bgr} is a useful adversarial input.
For its integer bracket table, 32 integer row identities with multipliers
$1,3,6$ force every derivation to map $L$ into $[L,L]$ in characteristic
$p>3$. Since it has class three, every derivation has cube zero. It thus
admits no positive integer grading, although finite-order automorphisms
can exist. The same argument extends to direct sums: off-diagonal
derivation blocks have central image, and the center is contained in
the derived algebra. Nevertheless the Rees cover of $m$ copies has
grades of dimensions $8m,4m,2m$ and the evaluation homomorphisms above.
The repository replays the integer identities and checks the public
cover maps, not a full quantum instrument on this input.
\end{remark}

\begin{lemma}[Copy-only lift]
Exactly $c$ fresh copies of $\rho_H$ suffice to produce
$\rho_{\widetilde H}$ on the public group $\exp(A)$.
\end{lemma}
\begin{proof}
Prepare the uniform address superposition and, on target $i$, apply
controlled $R_{\exp E_{s_i}(\log a)}$. Clean the public evaluation workspace
and discard all targets. Since $\tr(\rho_H R_g)=1_H(g)$, the retained
matrix entry at $a,b$ is
\[
 \frac1{|\exp A|}\prod_{i=1}^c
  1_H\bigl(\exp E_{s_i}(\log(b^{-1}a))\bigr)
 =\frac{1_{\widetilde H}(b^{-1}a)}{|\exp A|}.
\]
This is the desired coset density by Lemma~\ref{lem:homogeneous}.
No ordinary oracle is reconstructed from copies.
\end{proof}

\begin{proposition}[Refreshing correlated coset centers]
Let $P_a=|aH\rangle\langle aH|$. Suppose the supplied bank and side
information have the state
\[
 \Omega=\sum_{\mathbf a}\left(\bigotimes_{j=1}^M P_{a_j}\right)
               \otimes E_{\mathbf a},\qquad
 E_{\mathbf a}\succeq0,\quad\sum_{\mathbf a}\tr E_{\mathbf a}=1.
\]
Independent uniform left translations, with their randomization records
discarded, transform this bank into
$\rho_H^{\otimes M}\otimes\sum_{\mathbf a}E_{\mathbf a}$.
This uses public group sampling and multiplication, and no extra source
copies, source inverse, or knowledge of $H$.
\end{proposition}
\begin{proof}
Write $L_g|x\rangle=|gx\rangle$ and
$\mathcal T(S)=|G|^{-1}\sum_gL_gSL_g^\dagger$.
The transitive action on left cosets gives $\mathcal T(P_a)=\rho_H$
for every $a$. Tensor the channels and apply linearity to $\Omega$.
Alternatively,
\[
 \mathcal T(S)_{q,r}=|G|^{-1}\sum_z S_{zq,zr}
     =|G|^{-1}\tr(SR_{r^{-1}q}),
\]
so the existing uniform-address controlled-right-action and target-discard
channel implements the same refresh. The records of the random translations
are not retained as adaptive inputs to the receiver.
\end{proof}

\begin{remark}[Joint promise, not marginal tomography]
This is a standard twirling consequence, not a new quantum primitive.
The bank promise is sufficient, not necessary, and not certified from its
marginals. For $G=C_p$, odd $p$, $H=0$, let $|\chi_k\rangle$ be Fourier
states. The separable bank
\[
 \frac1p\sum_k |\chi_k\rangle\langle\chi_k|
       \otimes|\chi_{-k/2}\rangle\langle\chi_{-k/2}|
\]
has both marginals $I/p$ and is invariant under independent left twirls.
The two-target lift with evaluation points $1,2$ nevertheless yields
$|G\rangle\langle G|$, not $I/p$, because the character phases cancel.
The trace distance is $1-1/p$. Reusing a common translation for every
register, or conditioning on retained translation records, likewise does
not prove independent inputs. A joint source error $\delta$ and independent
uniform-sampling errors $\epsilon_j$ give output error at most
$\delta+\sum_j\epsilon_j$ by contractivity and a hybrid bound; marginal
errors cannot replace the joint error. No source-copy cost in the preceding
theorem is reduced: all consumed registers must actually be supplied.
\end{remark}

\section{The complete central instrument}
For this section let $G$ be any finite group, $H\le G$, and let
$Z\le Z(G)$ be elementary abelian over $\Fp$. Put $K=H\cap Z$.
For a central character $\chi_\lambda$, define
\[
 \Pi_\lambda=\frac1{|Z|}\sum_{z\in Z}\chi_\lambda(-z)R_z.
\]
Here $-$ denotes inverse in the additively written center.
The $\Pi_\lambda$ are orthogonal projectors summing to identity. They
can be measured with a uniform central register, controlled translations
and an abelian QFT. The original target remains the full group register.

\begin{lemma}[Normalized branch law]\label{lem:central}
The character outcome has probability $|K|/|Z|$ on $K^\perp$ and zero
otherwise. On every positive-mass outcome its normalized target state is
\[
 \rho_\lambda=\frac1{|G||K|}
       \sum_{h\in H,z\in Z}\chi_\lambda(-z)R_{hz}.
\]
For $g\in G$,
\[
 \tr(\rho_\lambda R_g)=
 \begin{cases}\chi_\lambda(z),&g=hz\in HZ,\\0,&g\notin HZ.\end{cases}
\]
The value is independent of the choice of decomposition $g=hz$.
\end{lemma}
\begin{proof}
Centrality makes $\rho_H$ commute with $\Pi_\lambda$.
The regular trace vanishes away from the identity. Thus
$\tr(\rho_H\Pi_\lambda)=|Z|^{-1}\sum_{z\in K}\chi_\lambda(-z)$,
which gives the probability. Dividing $\rho_H\Pi_\lambda$ by this
probability gives the displayed state. In its trace against $R_g$, a
nonzero term requires $g\in HZ$. Each such element has exactly $|K|$
decompositions, with the same character value because $\lambda$ kills
$K$. This multiplicity cancels the denominator $|K|$.
\end{proof}

Assume now that publicly computable automorphisms $\Delta_s$ preserve
$H$ and scale $Z$ by $s^d$ for every $s\in\Fp^\times$.
Take independent sources, measure their central labels
$\lambda_1,\ldots,\lambda_M$, and find a nonzero vector of scales with
\begin{equation}\label{eq:sde}
 \sum_{i=1}^M s_i^d\lambda_i=0.
\end{equation}
Zero-scale targets are consumed and discarded; $\Delta_0$ is never called.
Let $Q=G/Z$ and let $\sigma:Q\to G$ be any public section, not necessarily
a homomorphism. On a uniform $Q$ address, apply controlled
$R_{\Delta_{s_i}(\sigma(q))}$ to each active target and trace the targets.

\begin{proposition}[Pointwise quotient conversion]\label{prop:descent}
For every positive-mass tuple of central labels and every valid witness
to \eqref{eq:sde}, the output is exactly $\rho_{HZ/Z}$ on $Q$.
Neither $K$ nor a hidden coset representative is an operational input.
\end{proposition}
\begin{proof}
The output entry at $q,r$ is
\[
 \frac1{|Q|}\prod_{i:s_i\ne0}
 \tr\bigl(\rho_{\lambda_i}
 R_{\Delta_{s_i}(\sigma(r)^{-1}\sigma(q))}\bigr).
\]
Put $g=\sigma(r)^{-1}\sigma(q)$. If $g\notin HZ$, every active factor
is zero, since $\Delta_{s_i}(HZ)=HZ$; the active set is nonempty.
Otherwise write $g=hz$. Lemma~\ref{lem:central} makes the product
$\prod_i\chi_{\lambda_i}(s_i^d z)=1$ by \eqref{eq:sde}.
The remaining indicator is exactly coset equality for $HZ/Z$.
\end{proof}

This is an identity for complete measured branches, not postselection on
a favorable tuple. The output is independent of the retained classical
tuple. Fresh disjoint preparations therefore supply the independent
inputs required at the next stage. The proof does not authorize replacing
an IID supply by a correlated purification with the same marginals.

\section{Recovery and the published arithmetic dependency}
At stage $d=c,c-1,\ldots,2$, work in
$A/\bigoplus_{r>d}A_r$. Its grade $d$ is central with rank
$r_d=\dim\gamma_d(L)$, scales by $s^d$, and the hidden image remains
homogeneous. Proposition~\ref{prop:descent} applies successively.
The last quotient is $tL$ as an additive group, and the hidden image is
\emph{all} $t\mathfrak h$. Abelian Fourier samples reveal
$\mathfrak h^\perp$; a nullspace calculation returns $\mathfrak h$.
There is no hidden-intersection or normalizer recursion inside this core.

For fixed $d$, the diagonal-system theorem \cite[Sections 3.1--3.3]{is}
supplies a nonzero solution of \eqref{eq:sde} for arbitrary label columns
at a sufficient arity
\[
 M_d=B_d(r_d),\qquad
 B_d(r)=\left[d^{d(d-1)/2}(r+1)^d\right]^{\lceil\log_2(d+1)\rceil}.
\]
The displayed bound comes from the collision construction, not a heuristic
rank estimate. The exponent can first be reduced to $\gcd(d,p-1)$ and
converted back by Bezout. All returned witnesses are checked exactly.
There is no exhaustive search over field elements in the theorem claim.

\section{Full clocks and approximation errors}
Let $P=\prod_{d=2}^cM_d$, with empty product one, and take $m=n+t$
terminal samples, where $t=\lceil\log_p(3/\eps)\rceil$.
The rank failure probability is at most $1/((p-1)p^t)$.
The exact source and merge counts are
\[
 N_{\rm leaf}=mP,\qquad
 N_{\rm merge}=m\sum_{d=2}^c\prod_{j=2}^{d-1}M_j,
 \qquad N_{\rm read}=m.
\]
The original query or copy count is $cN_{\rm leaf}$.
Crucially, the full instrument count includes final Fourier readout:
\[
 N_{\rm inst}=N_{\rm leaf}+N_{\rm merge}+N_{\rm read}.
\]
Source preparation, central conversion and terminal QFT/measurement are
three separate error-bearing operations. Incorporate their internal
arithmetic and Fourier approximations into these full-channel allowances.
Here operational channel error means half the diamond norm distance;
for a fixed source preparation it means trace distance. Assign at most
$\eps/(3N_{\rm inst})$ error to each complete instrument.
A trace-distance hybrid bounds their total effect by $\eps/3$; no division
by small branch probabilities is used.

\begin{lemma}[Classical retries do not consume the source bank]
Suppose the verified classical SDE procedure has a known uniform expected
running-time upper bound $T$ on its legal public inputs. Capping each
fresh-seed attempt at $2T$ and trying
$\lceil\log_2(3N_{\rm merge}/\eps)\rceil$ times makes the abort probability
per merge at most $\eps/(3N_{\rm merge})$ using no additional source copies.
\end{lemma}
\begin{proof}
Markov gives failure probability at most $1/2$ per capped attempt.
Verification excludes accepting a false witness. Independent classical
seeds give the geometric abort bound even conditional on all prior quantum
records, since the input-dependent expected-time guarantee is uniform.
The labels have already been measured and stored, and no controlled
dilation is performed until a valid witness is found. Classical attempts
therefore reuse labels, not quantum copies. Include verification time in
$T$. If every attempt fails, abort the run with its charged probability.
\end{proof}

The case $N_{\rm merge}=0$ needs no SDE allowance. Uniform finite-field
preparations and QFTs have standard polynomial approximate implementations.
For arbitrary cyclic modulus one may use the constructions surveyed in
\cite[Chapters 3 and 5]{hales}; this does not assume an exact prime-modulus
QFT from a fixed finite gate set.
Reversible BCH, evaluation, dilations and right translations use public
fixed-degree arithmetic. The ideal schedule has polynomially many arithmetic
primitives at fixed $c$. With $\delta=\eps/(3N_{\rm inst})$, reversible
arithmetic and approximate Fourier circuits require polynomially many gates
in their input sizes and $\log(1/\delta)$. Thus precision and gate cost
remain polynomial in the input and $\log(1/\eps)$.
The rank, solver-abort and channel errors sum to at most
$\eps$. This proves the proposed explicit-Lie claim, subject to review of
the lemmas and imported algorithms.

The source-tree leaves need not all be stored simultaneously. With
depth-first preparation, at most $2+\sum_{d=2}^c(M_d-1)$ descent group
registers coexist when $c\ge2$; for $c=1$ only one is needed. This excludes
leaf-oracle workspaces, arithmetic ancillas and classical SDE storage,
which must also be charged. The product $P$ remains the time/copy cost.
No practical resource claim follows from these enormous polynomial bounds.

\section{Explicit input with small derived dimension}\label{sec:rank}
This is a proposed parameterized corollary of the same composition, not a
claim that arbitrary growing class is efficient or a new diagonal solver.
Continue to assume explicit genuine structure constants and $c<p$, but
now let $c$ vary. Write $r=\dim[L,L]$ and $r_d=\dim\gamma_d(L)$.

\begin{lemma}[Uniform public BCH arithmetic]
Given an $N$-dimensional Lie algebra of class at most $c<p$, its BCH product
can be evaluated with $O(c^4)$ Lie brackets and polynomially many additional
field operations, without enumerating free Lie words.
\end{lemma}
\begin{proof}
Set $Z(t)=\log(\exp X\exp(tY))$. The standard formal inverse-exponential
differential identity gives
\[
 Z'(t)=\sum_{k=0}^{c-1}b_k\operatorname{ad}_{Z(t)}^k(Y),\quad Z(0)=X,
 \qquad \sum b_kz^k=\frac{z}{1-e^{-z}}\pmod {z^c}.
\]
For an elementary derivation of this differential approach see
\cite[Section 4.1]{wehler}. Compute $b_k$ by reciprocal power series:
$f_k=(-1)^k/(k+1)!$, $b_0=1$ and
$b_k=-\sum_{i=1}^k f_i b_{k-i}$. Thus $b_1=+1/2$.
Given $Z_0,\ldots,Z_j$, polynomial bracket convolution finds the
coefficient of $t^j$ on the right, and division by $j+1$ gives $Z_{j+1}$.
Perform this for $0\le j<c$, then evaluate at $t=1$.
The universal class-$c$ BCH polynomial has $t$-degree at most $c$ and
denominators supported on primes at most $c$; the recurrence has the same
property. Reduction modulo $p$ is therefore valid and coefficientwise
uniqueness holds below degree $p$. No analytic convergence is required.
The direct convolution uses
$(c-1)c(c+1)(c+2)/6$ brackets; each dense bracket costs $O(N^3)$ field
operations. Reversible arithmetic and cleanup retain polynomial overhead.
\end{proof}

\begin{corollary}[Proposed derived-dimension parameter bound]
With the explicit input and source promises of Theorem~\ref{thm:main},
but allowing $c<p$ to vary, the proposed recovery has bounded-error time
\[
 2^{O(r^2\log(r+2))}\,
 \operatorname{poly}(n,\log p,\text{label/access width},\log(1/\eps)).
\]
The polynomial exponent is independent of $r$. For $r=0$ use abelian HSP.
\end{corollary}
\begin{proof}
After higher grades have been removed, choose a public complete flag of
the central grade $A_d$. Every flag prefix is an ideal invariant under
$\Delta_s$, since its action on this grade is scalar $s^d$. The hidden
image remains dilation invariant in each quotient. Apply
Proposition~\ref{prop:descent} separately to each one-dimensional factor;
the flag need not be adapted to the unknown hidden intersection.

The scalar diagonal theorem, \cite[Fact 2]{is}, needs $d+1$ labels and
time polynomial in $d,\log p$, uniformly in degree. Equivalently put
$\delta_d=\gcd(d,p-1)$ and use $\delta_d+1$ labels. A witness $x_i$ for
degree $\delta_d$ becomes $s_i=x_i^u$ on nonzero entries and $s_i=0$
otherwise, where $du\equiv\delta_d\pmod {p-1}$. This preserves the
equation and nonzeroness. Zero entries must stay zero even if $u=0$.
The fresh-source factor per terminal sample is consequently
\[
 P_p=\prod_{d=2}^c(\delta_d+1)^{r_d}
 \le \prod_{d=2}^c(d+1)^{r_d}.
\]
The nonzero lower-central terms strictly decrease, whence
$c\le r+1$, $r_d\le r-d+2$ and
\[
 D\le n+\tfrac12r(r+1),\qquad
 P_p\le (r+2)^{r(r+1)/2}.
\]
For $M=n+\lceil\log_2(3/\eps)\rceil$ final samples, there are $MP_p$
lifted leaves and $cMP_p$ original queries or copies. A forest with
arity at least two has at most $MP_p-M$ conversion nodes. Including leaf
and final-readout instruments, at most $2MP_p$ complete instruments need
error allowances. The earlier full-channel hybrid uses precision logarithmic
in this count. Public flag/quotient calculations and the preceding BCH
lemma have uniform polynomial cost in $D,c,\log p$. Combining these bounds
and the uniform scalar solver gives the claimed parameter dependence.
\end{proof}

If all $\delta_d\le b$, then $P_p\le(b+1)^{r(r+1)/2}$; for example this
holds with $b=2$ when the supplied prime is $p=2q+1$ with $q>c$ prime.
No safe-prime generation assumption is needed or claimed. The generic
sufficient polynomial regime is $r^2\log(r+2)=O(\log I)$ for explicit
input length $I$, and $r^2=O(\log I)$ in the bounded-index case.
Padding by abelian factors can manufacture such a regime without supplying
a meaningful application. Even the class-twelve profile $(13,11,10,\ldots,1)$
at $p=47$ costs over $3\cdot10^{28}$ original inputs for sixteen final
samples in this schedule; these are symbolic bounds, not observed runtime.

The earlier small-commutator algorithm \cite[Theorem 11]{ims} is polynomial
in input size plus $|G'|=p^r$; that particular estimate does not supply
the displayed uniform dependence on $\log p$. This is not a novelty proof
against all HSP literature. No growing-class opaque-input, general-exponent,
or natural hard application follows. The new finite arithmetic tests do not
independently validate the composition or implement a scalar solver/receiver.

\section{Direct descent for explicit prime-power Lie rings}\label{sec:ring}
This section gives an alternative parameter bound, not a strengthening of
the unrestricted-rank fixed-class outer theorem below. In particular it
does not assert an efficient solver for arbitrary diagonal systems over
torsion rings. All failures in the following randomized conversion are paid.

Suppose the input instead gives a genuine finite Lie ring $L$ with an
explicit additive decomposition $\bigoplus_{i=1}^n\mathbb Z/p^{a_i}$,
compatible bracket constants, and class $c<p$ for a supplied odd prime $p$.
The full $a_i\log p$ coordinate/order bit lengths belong to the input.
Lazard correspondence applies to these Lie rings
\cite{lazardring,lazardtorsion}; no
black-box chart acquisition is included here. Let $r$ be the minimum
additive generator count of $[L,L]$.

\begin{proposition}[Proposed direct ring parameter bound]
Under the same legal ordinary-oracle or coset-copy source contract, the
Rees argument has a direct bounded-error variant with time
\[
 [4(c+1)]^{(c-1)r}\,
 \operatorname{poly}(\text{explicit input bits},\log(1/\eps)).
\]
This is fixed-parameter time in $(c,r)$, not in $r$ alone.
\end{proposition}
\begin{proof}
Let $p^E$ annihilate $L$. Form
$A=\bigoplus_{d=1}^c t^d\gamma_d(L)$ with its actual finite-module
components, not free-module replacements. The evaluations $1,\ldots,c$
have Vandermonde determinant a unit modulo $p^E$. The earlier argument
uses only additive operations and unit scalars, so its hidden preimage
is still $\bigoplus t^d(\mathfrak h\cap\gamma_d(L))$.
Unit dilations preserve it. BCH identities and the uniform differential
recurrence are over $\mathbb Z[1/c!]$, whose denominators are units here.
Public submodules, sections and quotients use full integer relation lattices.

Remove the highest remaining grade one cyclic direct summand $Z$ at a
time. If $|Z|=p^a$ and $|H\cap Z|=p^b$, character orthogonality gives
independent uniform labels $\lambda\in p^b\mathbb Z/p^a$ on fresh sources.
The central contraction proof applies to finite abelian $Z$ verbatim.
At dilation degree $d$, consume $B=\gcd(d,p-1)+1\le d+1\le p$ targets.
If a label is zero, use scale one there and zero elsewhere. If all labels
are nonzero of equal valuation $v$, put $A_i=\lambda_i/p^v$ and solve the
scalar degree-$d$ equation modulo $p$ via \cite[Fact 2]{is} and the previous
power-index conversion. All $A_i$ are units. A nonzero witness coordinate
$s_j$ has derivative $dA_js_j^{d-1}$ a unit modulo $p$; keep other
coordinates fixed and Hensel-lift $s_j$ to modulus $p^{a-v}$. The resulting
active scales are units, inactive scales are zero, and
$\sum\lambda_i s_i^d=0\pmod {p^a}$. Every successful record therefore
produces the same correct quotient coset state.

Otherwise discard the entire batch and retry on fresh sources. If $b<a$,
the probability that all labels have valuation $b$ is
$(1-1/p)^B\ge1/4$; if $b=a$, acceptance is certain. Thus rejection costs
at most a factor four in expectation per cyclic conversion, independent
of $a$. This is a source-law argument, not a worst-case ring-SDE solver.

Write $r_d$ for the number of cyclic factors of $\gamma_d(L)$. A subgroup
of a finite abelian $p$-group with $r$ generators needs at most $r$
generators: $B[p]\subseteq A[p]$ and $d(B)=\dim_{\Fp}B[p]$ for $B\le A$.
Thus $\sum_{d=2}^c r_d\le(c-1)r$. Expected lifted leaves per
terminal sample are at most
$Q=\prod_{d=2}^c[4(\gcd(d,p-1)+1)]^{r_d}\le[4(c+1)]^{(c-1)r}$.
For $M=n+3$ final samples, cap the WHOLE run at $16MQ$ lifted leaves.
Markov bounds clock abort by $1/16$. Each successful or failed attempt
consumes at least two fresh targets, so all conversion forests together
have no more internal nodes than leaves. Including readout, at most
$32MQ+M$ complete instruments require precision allowances. Allocate
total channel error $1/16$ over this fixed cap; its precision overhead is
logarithmic. Every lifted leaf still costs $c$ original inputs.

The annihilator of the final hidden additive subgroup has at most $n$
generators. Its uniform samples generate it iff their images generate
its quotient by $p$, giving rank failure $<p^{-3}/(p-1)<1/16$.
One base run therefore succeeds with probability at least $13/16$.
Recover the hidden submodule and canonicalize its full preimage lattice
by Hermite normal form. Independent whole-run repetition and voting on
this canonical output amplify success with $O(\log(1/\eps))$ repetitions;
arbitrary Smith bases cannot be majority-voted. Joint input-bank error
is charged separately to the final $\eps$. All public arithmetic, the
scalar solver and Hensel lifting are polynomial in their bit sizes,
completing the proposed parameter bound.
\end{proof}

Two failures are important. First, in $\mathbb Z/9\times\mathbb Z/9$
with dilation degrees one and two and hidden subgroup $3\mathbb Z/9\times0$,
the nonunit scale $s=3$ cancels any degree-two central phase but falsely
returns the whole-quotient signal; phase cancellation alone is insufficient.
Second, the Lie ring over $\mathbb Z/5^4$ with $[x,y]=5y$ has derived
additive rank one and class four. The field inequality $c\le r+1$ cannot
be imported into this ring statement. Some mixed-valuation tuples have no
unit-scale solution, e.g. $(1,3,9)$ modulo $27$ at degree two.

The exact controls and expanded clock proof are in
\texttt{PRIME\_POWER\_REES\_DESCENT.md}. A module verifies and Hensel-lifts
supplied field witnesses; it does not implement their general scalar solver
or the quantum receiver. This alternative is not known to improve a useful
application, and for unrestricted derived rank at fixed class the earlier
outer reduction retains the better polynomial-time bound.

\section{Black-box extension: separate imported obligations}
This section is not an additional claim that an explicit Lie table is free.
Assume a uniquely encoded nilpotent black-box group, supplied generators,
coherent group operations, and fixed class. The standard presentation,
Sylow, subgroup, normalizer and quotient algorithms used in the nil-2 HSP
reduction are imported with their quantum preprocessing cost \cite{iss}.

For an exponent-$p$ local group with $p>c$, inverse-BCH group words give
Lazard addition and bracket. Their lengths depend only on $c$, with
denominators invertible modulo $p$. Enumerating generator brackets through
weight $c$ gives a polynomial-size additive spanning list. Abelian kernel
learning on its additive encoder produces a basis. Further kernel calls
recover bracket coordinates. All these calls use public group operations,
not the hidden source. The original generators alone need not span $L$.
Use bounded clocks, verified encodings, and amplified canonical kernel
outputs; checking a few relations does not prove kernel completeness.

There is an alternative local exponent-$p$ input conversion which avoids
learning the target's relations altogether. Let $F_{s,c}$ be the public
free class-$c$ nilpotent Lie algebra over $\mathbb F_p$ on the $s$ supplied
generators. Its dimension $d$ is at most $\sum_{j=1}^c s^j$, polynomial
only for fixed $c$. The universal property and Lazard correspondence give
a surjection $\pi:F_{s,c}\to\log G$. Evaluate it by public inverse-BCH
words on a nested-bracket basis. No target Lie coordinates or kernel
basis are needed. The map $\Phi=\exp\pi$ is generally noninjective;
$f\circ\Phi$ hides $K=\Phi^{-1}(H)$ with one original query. In the
copy-only model, controlled right translation by $\Phi(a)$ and target
discard give
\[
 \rho_{\rm out}[a,b]=|\exp F_{s,c}|^{-1}
  \operatorname{Tr}(\rho_H R_{\Phi(b)^{-1}\Phi(a)})
  =|\exp F_{s,c}|^{-1}1_K(b^{-1}a).
\]
This costs one original copy with no kernel-index penalty, provided the
public evaluation workspace is cleaned. Apply Theorem~\ref{thm:main}
to the explicit $F_{s,c}$ and project the recovered full hidden preimage
to $G$. Surjectivity, promised by the supplied generators, gives $H$.
An exact public word compiler now implements the general fixed-class
construction by homogeneous error correction, retaining the shorter
class-four formulas where valid. Certificates are checked in the free
associative algebra; no native quantum circuit is synthesized.

This route still needs Rees homogenization. For the Dixmier--Lister
target with free generators $a,b,c,d$ mapping to $e_1,e_2,e_3,e_4$,
the kernel contains $[a,d]+[b,[a,c]]$. A free degree dilation by $u$
maps its image to $(u^2-u^3)e_7$, which is nonzero at $u=2$, $p>3$.
Thus neither the kernel nor an arbitrary hidden preimage is homogeneous.
The free dimension can exceed the target encoding width $b$: use $d$
and $\dim\gamma_j(F_{s,c})$ in the core budgets, not the bounds by $b$
valid for the acquired-target-chart route. Each map evaluation also costs
$O(dW_c\log p)$ public multiplications and paid reversible cleanup.
The bridge does not remove general-exponent or small-prime reductions.

For copy transport into a public subgroup or quotient, one original copy
and controlled right translation along a forward section $\sigma$ yield
\[
 \rho_{\rm out}[q,r]=|Q|^{-1}
             1_H(\sigma(r)^{-1}\sigma(q)).
\]
Subgroup restriction has no index-sized postselection cost. A quotient
kernel must be known to lie in $H$ and be normal in the current public
subgroup. Its section need not be homomorphic. Flatten nested forward
sections and clean their workspaces; no source decoder is assumed.

\begin{corollary}[Proposed general nilpotent extension]
Combining Theorem~\ref{thm:main} with the preceding public input conversion
and the published nilpotent-group reductions gives a bounded-error
polynomial-time HSP algorithm for uniquely encoded groups of any fixed
nilpotency class, without a bound on prime factors or exponent. The source
may be an ordinary hiding oracle or fresh mixed coset copies.
\end{corollary}
\begin{proof}[Reduction details requiring separate review]
Class-one groups use the standard abelian HSP algorithm directly. For
$c\ge2$, every prime $p>c$ is odd, as required by the explicit-Lie statement.
The standard normalizer procedure reduces a Sylow problem to local hidden
groups of order at most $p$ \cite[Section 3]{iss}. For $p>c$ restrict such
a local call to the public subgroup of elements of order dividing $p$ and
use the explicit-Lie theorem. Do not restrict the original arbitrary hidden
group to its order-$p$ part just once.

For $p\le c$, the prime is bounded. Use the forward-only bounded-error
central-erasure/HSMC route \cite[Proposition 3]{ii}, not its distinct exact
theorem requiring a source inverse. A nonzero degree-$(p-1)$ diagonal
witness gives a nonempty binary zero-sum support by Fermat, so the cited
arithmetic theorem also provides a conservative bounded-prime subroutine.
The existing outer reduction then recovers the full hidden group.

One public step in that outer reduction requires an explicit qualification.
Maintain $H_0\le H$, $H_0\le K\le N_G(H_0)$ and put $B=K'H_0$.
To descend when $A=K/B$ is not elementary abelian, use the full preimage
of $A[p]$. Taking each arbitrary generator's top nontrivial $p$-power
does not suffice: $(1,0),(1,1)$ in $C_{p^2}\times C_p$ have identical
$p$-th powers. Instead acquire a complete integer row relation basis $R$
for $A=\mathbb Z^k/\operatorname{row}_{\mathbb Z}(R)$. Compute a basis of
\[
 C=\{a\in\mathbb F_p^k:aR=0\pmod p\},
 \qquad x_a=\frac{\widetilde a R}{p}\in\mathbb Z^k.
\]
The map $a\mapsto x_a$ modulo the relation lattice is an isomorphism
$C\cong A[p]$: surjectivity follows by writing $px=aR$, and injectivity
from the full rank of $R$. Evaluate the integer words $x_a$ in the
original quotient generators and adjoin every generator of $B$. This
generates precisely $L=\pi^{-1}(A[p])$, a proper subgroup exactly when
$A$ is not elementary. Do not reduce the integer words modulo $p$.
If $(H\cap K)/H_0$ is nontrivial, it has an order-$p$ element. Its image
in $K/B$ has order dividing $p$, so $H\cap L>H_0$, preserving the descent
invariant. Completeness of public relations is an amplified preprocessing
obligation, not certified merely by checking individual relations. These
computations add polynomial public work and no hidden-source reads.
Although $B$ is used for this public arithmetic, it need not lie in $H$;
it is not a legal source-quotient kernel. The source is restricted to $L$
and only quotiented by the known subgroup $H_0\le H$ as required.

If encoding length is $b$, polynomial clocks can cap local recovery calls
by $b(b+1)$ per branch across Sylows. Allocate failures to preprocessing,
chart acquisition, large-prime cores and small-prime calls. Conditional on
no earlier error, each adaptive input satisfies the next algorithm's
promise; a first-error union bound handles the entire execution. Bounded
clocks also cover erroneous histories. These are imported public algorithms
and a written composition, not implemented black-box routines in this repo.
\end{proof}

\section{Scope comparison and attempted refutations}
The classical group/Lie construction, interpolation, abelian HSP,
central Fourier measurement and diagonal-system solver are not new.
The nil-2 algorithm already uses automorphism-twisted right actions and
phase cancellation \cite{iss}. The bounded-prime fixed-class result
\cite{ii} does not give growing-prime polynomial time merely by treating
the prime as an input variable. The special semidirect groups in
\cite{is,morales} and quasi-Hamiltonian groups in \cite{morales} are
different scopes, not lower bounds on other algorithms.

For an explicit scope witness, consider $UT_9(\Fp)$ with $p>8$.
It has class eight, and with $u_{ij}=I+E_{ij}$, the derived elements
$a=u_{13},b=u_{35},c=u_{57},d=u_{79}$ satisfy
$[[a,b],[c,d]]=u_{19}\ne I$ (commutators $xyx^{-1}y^{-1}$).
Thus its derived subgroup is not abelian. It cannot be an abelian-by-cyclic
group of the special semidirect form. Also
$\langle u_{12}\rangle\langle u_{23}\rangle\ne
\langle u_{23}\rangle\langle u_{12}\rangle$, excluding quasi-Hamiltonian
scope. This is a structural non-subsumption check for those inspected
families, not a proof that no other prior algorithm covers this group.

The key falsifiers are precise. Removing an evaluation point can permit a
nonzero coefficient polynomial to vanish at all remaining points. Replacing
$\gamma_r(L)$ by unrestricted $L$ and dropping higher powers can destroy
evaluation homomorphisms. Replacing the cover by the associated graded
loses derived hidden directions. Nonnormal $H$ requires the stated right
action convention. Forbidden central outcomes cannot be normalized, and
all-zero scale witnesses erase the support test. Correlated source banks
can invalidate the tensor product calculation. Unrestricted growing class
defeats the fixed-class polynomial bound; Section~\ref{sec:rank} has a
separate derived-dimension restriction. None is a universal lower bound against other
approaches to HSP.

\section{What this draft does and does not establish}
The repository has finite source-to-answer controls for a genuine
class-three group with nonhomogeneous, nonnormal hidden lines, as well as
separate higher-class algebra and channel tests. Those tests cover specified
instances, not the universal theorem. The general SDE backend, quantum
receiver, black-box presentation acquisition and outer reduction are not
implemented end to end. A separate Lean draft is not compiler-verified.

Independent review should first try to refute
Lemma~\ref{lem:homogeneous} and Proposition~\ref{prop:descent} together,
then the black-box interface and resource composition. Novelty remains
unestablished. The growing-prime polynomial-phase application examined
separately has a direct known-technique quantum-learning baseline; it is
not evidence of a first efficient learner. No DCP, lattice, graph, code,
ECDSA, or Shor-level application is claimed. This document is intended to
make the actual proposed theorem easy to attack, not to promote passing
tests into a discovery.

\begin{thebibliography}{99}
\bibitem{bgr}
Y. Bahturin, M. Goze and E. Remm,
\emph{Group gradings on filiform Lie algebras},
\href{https://arxiv.org/abs/1308.2396}{arXiv:1308.2396}, Section 4.2.
\bibitem{lazard}
\emph{Model-theoretic properties of nilpotent groups and Lie algebras},
Section 2.3, \href{https://arxiv.org/abs/2310.17595}{arXiv:2310.17595}.
\bibitem{is}
G. Ivanyos and M. Santha,
\emph{Solving systems of diagonal polynomial equations over finite fields},
\href{https://arxiv.org/abs/1503.09016}{arXiv:1503.09016}, Sections 3--5.
\bibitem{iss}
G. Ivanyos, L. Sanselme and M. Santha,
\emph{An efficient quantum algorithm for the hidden subgroup problem in
nil-2 groups}, \href{https://arxiv.org/abs/0707.1260}{arXiv:0707.1260};
\href{https://www.irif.fr/~santha/Papers/iss12.pdf}{author's final version}.
\bibitem{ii}
M. Imran and G. Ivanyos,
\emph{Zero sum subsequences and hidden subgroups},
\href{https://arxiv.org/abs/2304.08376}{arXiv:2304.08376};
Quantum Information Processing 23, 14 (2024).
\bibitem{morales}
M. E. S. Morales,
\emph{The Hidden Subgroup Problem in Semidirect Products and
Quasi-Hamiltonian Groups},
\href{https://arxiv.org/abs/2608.05321}{arXiv:2608.05321} (2026).
\bibitem{hales}
L. R. Hales,
\emph{The Quantum Fourier Transform and Extensions of the Abelian Hidden
Subgroup Problem}, doctoral dissertation (2002),
\href{https://arxiv.org/abs/quant-ph/0212002}{arXiv:quant-ph/0212002}.
\bibitem{ims}
G. Ivanyos, F. Magniez and M. Santha,
\emph{Efficient quantum algorithms for some instances of the non-Abelian
hidden subgroup problem},
\href{https://arxiv.org/abs/quant-ph/0102014}{arXiv:quant-ph/0102014}, Theorem 11.
\bibitem{wehler}
J. Wehler, \emph{Lie Groups}, lecture notes, release 0.51 (2017), Section 4.1,
\href{https://www.mathematik.uni-muenchen.de/~wehler/LieAlgebrasLieGroupsScript.pdf}{author's notes}.
\bibitem{lazardring}
S. Cical\`o, W. A. de Graaf and M. Vaughan-Lee,
\emph{An effective version of the Lazard correspondence},
Journal of Algebra 352 (2012), 430--450,
\href{https://doi.org/10.1016/j.jalgebra.2011.11.031}{doi:10.1016/j.jalgebra.2011.11.031}.
\bibitem{lazardtorsion}
O. Garaialde Oca\~na and J. Gonz\'alez-S\'anchez,
\emph{Transporting cohomology in Lazard correspondence},
\href{https://arxiv.org/abs/1405.4654}{arXiv:1405.4654}, Theorem 1.
\end{thebibliography}
\end{document}
